\documentclass[11pt]{article}

\usepackage[T1]{fontenc}
\usepackage{lmodern}
\usepackage[margin=1in]{geometry}
\usepackage{microtype}
\usepackage{amsmath,amssymb,amsthm}
\usepackage{booktabs,array}
\usepackage{enumitem}
\usepackage{needspace}
\usepackage{xurl}
\usepackage[hidelinks]{hyperref}

\newtheorem{theorem}{Theorem}
\newtheorem{lemma}[theorem]{Lemma}
\newtheorem{proposition}[theorem]{Proposition}
\newtheorem{corollary}[theorem]{Corollary}
\newcommand{\Q}{Q}
\newcommand{\interior}{\operatorname{int}}
\newcommand{\R}{\mathbb R}
\newcommand{\ind}{\mathbf 1}
\newcommand{\abs}[1]{\left|#1\right|}
\newcolumntype{L}[1]{>{\raggedright\arraybackslash}p{#1}}

\title{A Parent-Aware Computer-Assisted Lower Bound for\\Packing 17 Unit Squares in a Square}
\author{Mira\\\small \url{https://github.com/Mira-acc/17squares}}
\date{September 2026}
\hypersetup{
  pdftitle={A Parent-Aware Computer-Assisted Lower Bound for Packing 17 Unit Squares in a Square},
  pdfauthor={Mira},
  pdfsubject={Computer-assisted discrete geometry},
  pdfkeywords={square packing, weighted covering, exact arithmetic, parent-aware certificates}
}
\emergencystretch=2em
\raggedbottom

\begin{document}
\maketitle

\begin{abstract}
Let $s(n)$ be the least side length of a square containing $n$ unit squares
with pairwise disjoint interiors and arbitrary orientations. We prove
\[
  s(17)>\frac{46129999999859}{9997499999900}
  =4.614153538416645696808\ldots.
\]
The proof assigns a finite nonnegative measure to the container and selects a
smaller closed square inside each possible packed square. An exact finite
calculation shows that every selected square captures more than one seventeenth
of the total mass. Following the parent-aware reduction of Guzhou0806, we
check only centers attainable by the enclosing square. We then sharpen the
bound by moving the selected squares within their parents while reusing the
verified coverage. As a separate result, we construct intersection-trigger
charges and give an exact 29-square example that separates them from
nonnegative spatial measures.
\end{abstract}

\section{Introduction}

Write $\Q_L=[-L/2,L/2]^2$ for $L>0$. For a positive integer $n$, let $s(n)$
be the minimum $L$ for which $\Q_L$ contains $n$ unit squares with pairwise
disjoint interiors. Boundary contacts
are allowed, and each square may rotate independently. The minimum is attained;
a compactness argument is given below. For $n=17$, Bidwell's construction gives
\[
  s(17)\le L_B,\qquad L_B=4.6755300936045509516\ldots,
\]
as recorded in Friedman's survey and Ellsworth's
catalogue~\cite{Friedman2009,Ellsworth}. We establish the following lower bound.

\begin{theorem}\label{thm:main}
There is no packing of seventeen unit squares in a square of side
\[
  S_* = \frac{46129999999859}{9997499999900}.
\]
Consequently, $s(17)>S_*$.
\end{theorem}

The proof uses a finite nonnegative measure. Inside every possible packed
square, called a \emph{parent}, we select a smaller closed square, called a
\emph{core}. The cores lie strictly inside their parents, so cores selected
from a packing are disjoint even when the parents touch. We arrange that
every selected core captures mass at least $\gamma$, while the total mass is
less than $17\gamma$. Seventeen parents would therefore require more mass
than is available.

The difficulty is to prove the mass bound for every translation and
orientation of a parent. A finite catalogue assigns one core shape to each
interval of parent orientations. For that interval we check all centers where
a parent can occur. This domain can be smaller than the set of centers where
the core alone fits, and thus imposes fewer coverage constraints.

\paragraph{Contributions and dependencies.}
Weighted covering extends the classical unavoidable-point argument; recent
constructions and exact verification methods include those of Burns,
Massaccesi, and Levy~\cite{Burns,Massaccesi,Levy}. We use Guzhou0806's R012
parent-angle catalogue and legal-center reduction~\cite{Guzhou}, whose
antecedents include Levy's parent-center restrictions. Our base certificate
refines the angle catalogue and changes both the weights and the spatial
support, giving
\[
  s(17)>\frac{18452}{3999}=4.61415353838459614903\ldots.
\]
A recentering argument then proves Theorem~\ref{thm:main}, improving this base
constant by approximately $3.205\times10^{-11}$ without another translation
sweep. The intersection-trigger result is independent of this numerical bound:
it shows a strict advantage over spatial measures on a specified finite family,
not coverage of every possible parent.

Sections~\ref{sec:certificate}--\ref{sec:verification} give the certificate
criterion, the complete finite reduction, and the base computation.
Section~\ref{sec:recenter} proves the sharpening, and Section~\ref{sec:triggers}
treats triggers. Reproduction instructions and supplementary results, including
the dated history of improvements, are in the appendices. The computational
claims use rational data and exact arithmetic; the geometric reductions are
proved here. The complete argument has not been formalized in a proof assistant.

\section{Parent-aware certificates}\label{sec:certificate}

\subsection{Counting disjoint cores}

A \emph{charge} assigns a nonnegative number $C(K)$ to each closed core $K$.
Its \emph{budget} is any number $\mathcal B$ such that
\[
  \sum_i C(K_i)\le\mathcal B
\]
for every finite pairwise-disjoint family of cores. For a finite nonnegative
measure $\mu$, the charge $C(K)=\mu(K)$ has budget $M=\mu(\R^2)$.

\begin{lemma}[Disjoint-core counting]\label{lem:budget}
Suppose every admissible parent admits a closed core strictly inside its
interior with charge at least $\gamma>0$. If the charge has budget
$\mathcal B<n\gamma$, then $n$ admissible parents cannot have pairwise
disjoint interiors.
\end{lemma}

\begin{proof}
Choose one such core from each parent. These cores are disjoint, so their
total charge is both at least $n\gamma$ and at most $\mathcal B$, a
contradiction.
\end{proof}

The strict containment is essential for atomic measures: an atom on a shared
parent boundary must not be counted twice. The numerical certificate below
uses only point masses. The more general language will also apply to the
charges in Section~\ref{sec:triggers}.

\subsection{Orientations, containment, and legal centers}

Let $D_4$ be the eight symmetries of the container. We choose a $D_4$-invariant
measure, so an individual parent may be reflected into orientation
$\phi\in[0,\pi/4]$ without changing any core mass. The square's own quarter-turn
symmetry accounts for the remaining orientations. This reduction imposes no
symmetry on a hypothetical packing.

Use the half-angle parameter $u=\tan(\phi/2)$, whose physical folded range is
$[0,\sqrt2-1]$. For the rational computations it is convenient to define,
more generally for $0\le u<1$,
\begin{equation}\label{eq:trig}
  c(u)=\frac{1-u^2}{1+u^2},\qquad
  s(u)=\frac{2u}{1+u^2},\qquad
  f(u)=c(u)+s(u).
\end{equation}
The first two quantities are the cosine and sine of $2\arctan u$; $f(u)$ is
the axis-parallel width of a unit square in that direction.

A catalogue row consists of an interval $I_j=[\alpha_j,\beta_j]\subset[0,1)$,
a core parameter $t_j\in[0,1)$, and a core side $B_j>0$. The parent side is
$A>0$. To compare the two directions, define
\begin{equation}\label{eq:g}
  g(t,u)=c(t)c(u)+s(t)s(u)+\abs{s(t)c(u)-c(t)s(u)}.
\end{equation}
When the relative angle is at most $\pi/4$ in absolute value, $B_jg(t_j,u)$
is the width of the core across either parent normal.

\begin{lemma}[Interval containment]\label{lem:contain}
Suppose the angles represented by $t_j$ and each endpoint of $I_j$ differ
by at most $\pi/4$ in absolute value. Put
\[
  G_j=\max\{g(t_j,\alpha_j),g(t_j,\beta_j)\}.
\]
If $B_jG_j<A$, then the concentric closed core lies strictly inside the
side-$A$ parent for every $u\in I_j$.
\end{lemma}

\begin{proof}
The largest absolute angular difference from a fixed direction occurs at an
interval endpoint. On $[0,\pi/4]$, the function $\cos v+\sin v$ is increasing.
Thus the core width across either parent normal is at most $B_jG_j<A$.
Both normal-coordinate inequalities are strict, proving the claim.
\end{proof}

The relative-angle hypothesis is checked rationally at each endpoint: the dot
product of the direction vectors must be positive and at least the absolute
value of their cross product. No sampling assertion is used between endpoints.

A side-$A$ parent with parameter $u$ fits in $\Q_L$ precisely when its center
lies in the square of half-side $(L-Af(u))/2$, provided this is nonnegative.
The complete envelope for an interval has a simple exact description.

\begin{lemma}[Legal-center envelope]\label{lem:domain}
Set
\[
  m_j=\min\{f(\alpha_j),f(\beta_j)\},\qquad
  d_j=\frac{L-Am_j}{2}.
\]
If $d_j\ge0$, the union of the legal center sets over $I_j$ is exactly
$D_j=[-d_j,d_j]^2$.
\end{lemma}

\begin{proof}
Since
\[
  f'(u)=\frac{2(1-2u-u^2)}{(1+u^2)^2},
\]
$f$ has a single interior maximum and no interior minimum on $[0,1)$.
Its minimum on $I_j$ is therefore $m_j$. The nonempty center squares are
nested, and the endpoint attaining $m_j$ supplies the largest one.
\end{proof}

In particular, every point of $D_j$ is attained by a legal parent at an
endpoint orientation. When an interval extends slightly beyond the folded
physical range, using its full envelope merely checks additional centers.

\Needspace{17\baselineskip}
\subsection{The certificate criterion}

\begin{theorem}[Parent-aware certificate]\label{thm:certificate}
Fix a positive integer $n$ and $L,A>0$. Let $\mu$ be a finite nonnegative
$D_4$-invariant measure supported in $\Q_L$, and let $\gamma>0$. Suppose a
finite catalogue has the following properties:
\begin{enumerate}[label=(\roman*),leftmargin=2.3em]
\item Its intervals cover $[0,\sqrt2-1]$.
\item Every row satisfies Lemma~\ref{lem:contain}, and its domain $D_j$ from
Lemma~\ref{lem:domain} has $d_j>0$.
\item Every core from row $j$ centered in $D_j$ has mass at least $\gamma$.
\item $M:=\mu(\R^2)<n\gamma$.
\end{enumerate}
Then $n$ side-$A$ squares cannot pack in $\Q_L$, and $s(n)>L/A$.
\end{theorem}

\begin{proof}
Fold each parent separately and choose a catalogue interval containing its
orientation. Its center lies in $D_j$, so the specified concentric core has
mass at least $\gamma$ and lies strictly inside the parent. Apply the inverse
reflection to the core. Invariance of $\mu$ preserves its mass, and strict
containment is preserved. Lemma~\ref{lem:budget} excludes a packing.

Scaling by $A$ excludes a unit-square packing at side $L/A$. To obtain the
strict inequality, note that the minimum defining $s(n)$ is attained. A grid
packing supplies a finite upper bound on the container side; centers and
orientations then range over compact sets. Containment is closed under limits.
Interior overlap is an open condition on a pair of square poses, so
interior nonoverlap is closed. A sequence of packings whose container sides
approach the infimum therefore has a feasible limit. Since the endpoint
$L/A$ has been excluded, that minimum is strictly larger.
\end{proof}

The remaining universal condition is (iii). We next reduce it to a finite
arrangement calculation and prove that the calculation includes all boundary
placements.

\section{Finite verification and the base bound}\label{sec:verification}

\subsection{Capture rectangles}

For the computation, write $\mu=\sum_{k=1}^N\omega_k\delta_{p_k}$ with
$\omega_k\ge0$. Assume $L,A$, the catalogue parameters, the atom coordinates,
and the weights are rational.
Fix one row and abbreviate its parameters by $t,B,d$. Write $c=c(t)>0$ and
$s=s(t)\ge0$. Rotate a potential core center $(x,y)$ into the core frame:
\[
  \xi=cx+sy,\qquad \zeta=-sx+cy.
\]
For a spatial atom $p$, let $(\xi_p,\zeta_p)$ be its rotated coordinates.
The closed core captures $p$ exactly when its center belongs to
\begin{equation}\label{eq:capture}
  R_p=[\xi_p-B/2,\xi_p+B/2]\times
      [\zeta_p-B/2,\zeta_p+B/2].
\end{equation}
Thus the mass at a center is the sum of weights of the rectangles containing
it. The legal-center square becomes the rotated square
\begin{equation}\label{eq:rotated-domain}
  \Omega=\{(\xi,\zeta):\abs{c\xi-s\zeta}\le d,
                             \abs{s\xi+c\zeta}\le d\}.
\end{equation}
All rectangle edges and domain vertices are rational.

\subsection{Which cells must be checked?}

Split the $\xi$-axis at every vertical edge of a capture rectangle and at the
domain extremes $\pm d(c+s)$. Between consecutive events, the rectangles
active in the $\xi$ direction are fixed. Split the $\zeta$-axis at all
horizontal rectangle edges. Each resulting open rectangular cell has a
constant captured mass. We must identify exactly those cells meeting the
interior of $\Omega$.

For an open vertical slab $(a,b)$ inside the domain's horizontal extent,
consider the projection of
$\Omega\cap([a,b]\times\R)$ onto the $\zeta$-axis. Convexity makes this a
closed interval $[\ell,h]$. The projection of the corresponding open
interior intersection is $(\ell,h)$: the closed intersection is the closure
of its interior, and linear projection of that nonempty convex interior is
an open interval. Consequently a horizontal cell $(v_k,v_{k+1})$ is reachable
with positive area exactly when
\begin{equation}\label{eq:reachable}
  \max\{v_k,\ell\}<\min\{v_{k+1},h\}.
\end{equation}
It is the projection of the \emph{whole slab} that matters; a cross-section
at its midpoint can miss reachable cells.

Here is an exact formula for that projection. If $s>0$, the lower and upper
boundaries of $\Omega$ are
\begin{align*}
  l(\xi)&=\max\{(c\xi-d)/s,\,(-d-s\xi)/c\},\\
  h_0(\xi)&=\min\{(c\xi+d)/s,\,(d-s\xi)/c\}.
\end{align*}
The convex function $l$ is minimized at $d(c-s)$, and the concave function
$h_0$ is maximized at $-d(c-s)$. Define
$\operatorname{clamp}(z;a,b)=\min\{b,\max\{a,z\}\}$. Then
\begin{equation}\label{eq:projection}
  \ell=l(\operatorname{clamp}(d(c-s);a,b)),\qquad
  h=h_0(\operatorname{clamp}(-d(c-s);a,b)).
\end{equation}
When $s=0$, the domain is axis aligned and $[\ell,h]=[-d,d]$.
These formulas also handle a domain vertex lying inside the slab.

The finite algorithm sums the active weights on every horizontal cell
satisfying~\eqref{eq:reachable}, for every vertical slab. Its smallest sum is
the smallest mass attained away from the arrangement and domain boundaries.
All comparisons use rational arithmetic.

\begin{lemma}[Boundary completeness]\label{lem:boundary}
For nonnegative atomic weights and $d>0$, the minimum over the reachable open
cells equals the minimum over the whole closed domain $\Omega$.
\end{lemma}

\begin{proof}
Every point of $\Omega$ is a limit of interior points avoiding the finitely
many rectangle edges. A subsequence lies in one fixed open cell. Every
capture rectangle containing that cell is closed, so it also contains the
limit point. Nonnegative weights imply that the mass at the limit is at
least the mass in that cell. Conversely, each reachable open cell contains
a point of $\Omega$. The two minima are therefore equal.
\end{proof}

This proves both obligations of the finite reduction: every tested mass is
realized by a legal center, and no untested placement has smaller mass. The
implementation additionally reconstructs a rational center for each reported
minimum and recounts all atoms there. That recount checks the reported witness;
the arrangement argument above establishes exhaustiveness.

\Needspace{14\baselineskip}
\subsection{The verified certificate}

\begin{proposition}[Computational base certificate]\label{prop:base}
For
\[
  L=\frac{4613}{1000},\qquad A=\frac{3999}{4000},
\]
the accompanying rational data satisfy Theorem~\ref{thm:certificate} with
$n=17$ and
\[
  M=\frac{4249997911}{250000000}=16.999991644,
  \qquad
  \gamma=\frac{100000003}{100000000}=1.000000030.
\]
In particular, $s(17)>18452/3999$.
\end{proposition}

\begin{proof}[Computer-assisted verification]
The measure has 3,280 atoms in 417 positive $D_4$-orbits. Its 4,391 catalogue
intervals meet exactly, start at zero, and end at the rational value
$207107/500000>\sqrt2-1$. Each row satisfies the endpoint containment checks
and has $d_j>0$. The exact sweep just described verifies mass at least
$\gamma$ throughout every $D_j$. Finally,
\[
  17\gamma-M=\frac{4433}{500000000}>0.
\]
Theorem~\ref{thm:certificate} gives the conclusion. Appendix~\ref{app:replay}
identifies the data, replay procedure, and arithmetic assumptions.
\end{proof}

\section{Reusing coverage by recentering cores}\label{sec:recenter}

Concentricity helped construct the base certificate, but the counting proof
requires only strict containment. We can therefore shrink the parents and
container slightly while moving each chosen core back into a center domain
where its mass has already been verified. The measure stays fixed in the
common container-centered coordinate system.

\subsection{A general transfer principle}

Let $K_j$ denote a core shape centered at the origin, and suppose every
translate $y+K_j$ with $y\in D_j$ has charge at least $\gamma$. For a parent
shape $P$ centered at the origin, define
\[
  E_j(P)=\{v:v+K_j\subset\interior P\}.
\]
A parent centered at $x$ can use entry $j$ exactly when some $y\in D_j$
satisfies $y-x\in E_j(P)$. Equivalently, $x\in D_j-E_j(P)$, where
$D-E=\{y-v:y\in D,\ v\in E\}$.
For centrally symmetric $P$ and $K_j$, the set $E_j(P)$ is also centrally
symmetric, so $D_j-E_j(P)=D_j+E_j(P)$.

\begin{proposition}[Recentered-core transfer]\label{prop:recenter}
Suppose $C$ has budget $\mathcal B<n\gamma$. If every admissible parent
$P+x$ satisfies
\[
  x\in\bigcup_j(D_j-E_j(P)),
\]
then $n$ such parents cannot have pairwise disjoint interiors.
\end{proposition}

\begin{proof}
Membership supplies an index $j$ and a certified center $y$ with
$y+K_j\subset\interior(P+x)$. The selected core has charge at least
$\gamma$. Apply Lemma~\ref{lem:budget}.
\end{proof}

No continuity of the atomic mass under translation is assumed: the mass
statement is used only at an already verified center.

\subsection{An exact criterion for square domains}

Let a new parent have side $a>0$, orientation $\phi\in[0,\pi/4]$, and
container side $H>0$. Define
\[
  F(\phi)=\cos\phi+\sin\phi,\qquad
  q=\frac{H-aF(\phi)}2.
\]
Thus its legal center set is $[-q,q]^2$ when $q\ge0$, and
$f(u)=F(2\arctan u)$ in the earlier notation. Suppose the old core has side
$B$, angle $\theta$, and certified center domain $[-d,d]^2$. Its width
across either parent normal is
\[
  w(\phi)=B\bigl(\abs{\cos(\theta-\phi)}+
                         \abs{\sin(\theta-\phi)}\bigr).
\]
Write $[z]_+=\max\{z,0\}$.

\begin{lemma}[Square recentering]\label{lem:recenter}
If $q\ge0$ and $d\ge0$, every legal parent center admits a strictly interior
core with center in $[-d,d]^2$ if and only if
\begin{equation}\label{eq:recenter}
  w(\phi)+F(\phi)[H-aF(\phi)-2d]_+<a.
\end{equation}
\end{lemma}

\begin{proof}
Clamp each coordinate of a parent center to $[-d,d]$. Each coordinate moves
by at most $[q-d]_+$, so displacement along either parent normal is at most
$F(\phi)[q-d]_+$. Adding the core half-width proves sufficiency.

For necessity when $q>d$, choose the parent center $(q,q)$. Every available
core center has both coordinates at most $d$, so displacement along the
normal $(\cos\phi,\sin\phi)$ has magnitude at least $F(\phi)(q-d)$.
Strict containment forces~\eqref{eq:recenter}. When $q\le d$, concentric
placement is available, and the necessary and sufficient condition is
$w(\phi)<a$.
\end{proof}

\subsection{The scalar sharpening}

For row $j$ of the base certificate, retain $m_j,G_j,d_j$ and set
\[
  \eta_j=A-B_jG_j>0.
\]
Choose $0<\varepsilon<A$ and $0<\delta<L$, and put
$a=A-\varepsilon$, $H=L-\delta$. Let $F_j$ be an upper bound for $f(u)$ on
$I_j$. For any folded physical parameter $u\in I_j$,
\begin{align*}
  H-af(u)-2d_j
    &=A(m_j-f(u))+\varepsilon f(u)-\delta\\
    &\le\varepsilon F_j-\delta.
\end{align*}
Moreover, $w(2\arctan u)=B_jg(t_j,u)\le A-\eta_j$.
Substitution into~\eqref{eq:recenter} shows that the sufficient condition
\begin{equation}\label{eq:scalar}
  \eta_j>\varepsilon+F_j[\varepsilon F_j-\delta]_+
\end{equation}
allows every legal new parent in that interval to use its old core and center
domain. Empty legal-center sets impose no obligation.

\begin{proof}[Proof of Theorem~\ref{thm:main}]
Take
\[
  \varepsilon=10^{-11},\qquad \delta=\frac{141}{10^{13}}.
\]
For each interval not containing $\sqrt2-1$, use the larger endpoint value
of $f$ as $F_j$. For an interval containing that maximum, use the rational
upper bound $1414213563/10^9>\sqrt2$. Exact evaluation of all 4,391
inequalities~\eqref{eq:scalar} gives
\[
  \min_j\bigl(\eta_j-\varepsilon
       -F_j[\varepsilon F_j-\delta]_+\bigr)>10^{-15}.
\]
The finite arithmetic check is included in the reproduction procedure.

Each new parent therefore admits a core whose mass is at least the same
$\gamma$ as in Proposition~\ref{prop:base}. Fold and unfold each parent as
before; both containers have the same symmetry center, and the old measure
remains $D_4$ invariant. The old total mass $M$ is still a valid budget,
even if some of its atoms lie outside the smaller container. Since
$M<17\gamma$, Proposition~\ref{prop:recenter} excludes seventeen side-$a$
parents in $\Q_H$. Scaling and the compactness argument give
\[
  s(17)>\frac{H}{a}
  =\frac{L-\delta}{A-\varepsilon}
  =\frac{46129999999859}{9997499999900}.
\]
\end{proof}

\section{Intersection-trigger charges}\label{sec:triggers}

The preceding bound uses additive spatial mass. We now construct charges that
also reflect incompatibility between possible square poses. This section gives
a general budget rule and a finite separation result; a packing improvement
would additionally require coverage for every legal parent.

\subsection{Intersecting triggers and their budget}

Let $T_1,\ldots,T_r$ be nonempty finite sets of points, called triggers, with
$T_i\cap T_j\ne\varnothing$ for $i\ne j$. Define
\[
  C_T(K)=\ind\{T_i\subseteq K\text{ for some }i\}.
\]

\begin{proposition}\label{prop:trigger}
The charge $C_T$ has budget one on pairwise-disjoint cores.
\end{proposition}

\begin{proof}
If two cores activate the same trigger, they share a point because that
trigger is nonempty. If they activate different triggers, they share a point
in their intersection. Thus two disjoint cores cannot both be charged.
\end{proof}

For example, Levy's two-of-three threshold charge~\cite{Levy} uses the three
two-element subsets of a three-point set as triggers. More generally, for a
finite set $X$, integer $k\ge1$, and weight $w\ge0$, a threshold charge
activates when a core contains at least $k$ points of $X$. Its budget is
$w\lfloor |X|/k\rfloor$, since disjoint cores have disjoint traces on $X$.
For arbitrary nonempty triggers, their matching number---the largest number
of pairwise-disjoint triggers---is a valid budget.

To obtain triggers from geometry, start with square poses
$V_1,\ldots,V_r$ having pairwise-overlapping interiors, with $r\ge2$. Choose
\[
  p_{ij}=p_{ji}\in\interior V_i\cap\interior V_j,
  \qquad T_i=\{p_{ij}:j\ne i\}.
\]
Then $T_i\subset\interior V_i$ and $p_{ij}\in T_i\cap T_j$. The resulting
budget-one charge activates on every $V_i$. Because the finitely many
witnesses lie strictly inside their squares, activation persists for a
sufficiently small change in each pose.

\subsection{A finite separation from spatial measures}

\begin{proposition}[29-core separation]\label{prop:separation}
There are 29 closed squares $K_i$ of side $499/500$ and an intersection-trigger
charge with budget one such that $C_T(K_i)=1$ for every $i$, whereas every
nonnegative Borel measure $\nu$ satisfying $\nu(K_i)\ge1$ for all $i$ has
\[
  \nu(\R^2)\ge\frac{81407286808}{59431493389}
  =1.3697668048682659\ldots>1.
\]
\end{proposition}

\begin{proof}[Computer-assisted construction and proof]
The accompanying data specify 29 rational unit-square poses $V_i$, positive
rational weights $\lambda_i$, and 406 pair witnesses. Every required witness
has inward normal distance greater than $1/1000$ from the boundary of its
unit square. It therefore lies inside the concentric closed square $K_i$ of
side $499/500$. The 29 triggers each contain 28 witnesses, and their pairwise
intersections prove $C_T(K_i)=1$ and budget one.

The exact weighted sum and maximum pointwise load are
\[
  W:=\sum_i\lambda_i=\frac{81407286808}{79999999999},\qquad
  \Lambda:=\max_{x\in\R^2}\sum_i\lambda_i
             \ind_{\interior V_i}(x)=\frac{59431493389}{79999999999}.
\]
To check the latter equality, partition the plane at all vertex abscissae and
all edge-intersection abscissae. In each open vertical slab, edge order is
fixed, so a rational cross-section enumerates every open cell's covering
set. Coincident events are grouped. At any boundary point, all open squares
containing that point also contain a neighborhood of it; hence no boundary
load exceeds the open-cell maximum. A maximizing witness is recounted
directly. The rational data and replay are specified in
Appendix~\ref{app:replay}.

Since $K_i\subset\interior V_i$, the same upper bound $\Lambda$ holds for
the weighted indicators of the closed cores, including their boundaries.
For any finite measure $\nu$ as in the statement,
\[
  W\le\sum_i\lambda_i\nu(K_i)
    =\int\sum_i\lambda_i\ind_{K_i}\,d\nu
    \le\Lambda\nu(\R^2).
\]
Division by $\Lambda$ proves the stated bound. Infinite total mass satisfies
it trivially.
\end{proof}

\subsection{Coverage geometry and symmetry}

At a fixed core direction, let $(\xi_p,\zeta_p)$ be the coordinates of a
trigger point in the core frame. The centers of side-$B$ cores containing
all of $T_i$ form the rectangle
\[
  R_i=\left[\max_{p\in T_i}\xi_p-B/2,
             \min_{p\in T_i}\xi_p+B/2\right]
       \times
       \left[\max_{p\in T_i}\zeta_p-B/2,
             \min_{p\in T_i}\zeta_p+B/2\right],
\]
where a reversed interval means the rectangle is empty. Thus $C_T$ is the
indicator of $\bigcup_{i=1}^rR_i$. A finite vertical-slab decomposition again
gives exact coverage cells. Rectangles belonging to a single trigger charge
must be evaluated as a union: an overlapping pair contributes one, not two.

One can also start from an overlap graph $G$, choosing a shared interior
witness for each edge and a private interior witness for each isolated vertex.
Each vertex trigger contains its incident edge witnesses, or its private
witness if the vertex is isolated. Disjoint charged
cores must select distinct, nonadjacent vertices, so any proved upper bound
$b$ on the independence number of $G$ gives a budget $b$.

Nonnegative sums of charges have the sum of their budgets, with no
independence assumption. In particular,
\[
  C(K)=\mu(K)+\sum_h w_h C_{T^{(h)}}(K),\qquad
  \mathcal B=M+\sum_h w_h b_h
\]
satisfies Lemma~\ref{lem:budget} for $w_h\ge0$. The recentering proposition
also applies once the necessary coverage is proved. To use the folded-angle
certificate, however, the \emph{total charge} must be $D_4$ invariant.
Otherwise coverage must be checked over the full orientation range. Averaging
$C(gK)$ over $g\in D_4$ preserves its budget and gives an invariant charge,
but its coverage must be established for that averaged charge. The finite
separation above does not use an orientation fold.

\paragraph{Acknowledgments.}
The direct mathematical and software dependencies are credited above and in
the artifact documentation. The research, computation, and exposition were
developed with substantial AI assistance under human direction. This is a
computer-assisted working preprint; independent external review remains open.

\appendix
\section{Reproducibility and arithmetic assumptions}\label{app:replay}

The public artifact~\cite{MiraRepo} contains the exact inputs, checkers, and
result records. The paper uses centered coordinates; data stored in
$[0,L]^2$ are translated by $(-L/2,-L/2)$ when interpreted in this notation.
This common change of origin also applies to the smaller container in the
recentering argument.

\subsection{Packing certificate}

The directory \path{certificates/lower_bound_4p614153/} contains
\path{certificate.json}, the authoritative base input; \path{verify.py} and
\path{parent_sweep.cpp}, the exact replay; and \path{sharpening.py}, the finite
recentering check. An orbit row $(x,y,w)$ expands into all distinct $D_4$
images of $(x/10^5,y/10^5)$, each with weight $w/D$, where $D$ is the recorded
weight denominator. A catalogue row stores the four rational values
$(\alpha_j,\beta_j,t_j,B_j)$.

The loader rejects duplicate orbit sites, negative weights, inconsistent
mass totals, incomplete interval coverage, failed strict-containment checks,
and degenerate center domains. It also verifies that the cores fit throughout
the supplied domains. The default input's exact byte identity is checked
against the SHA-256 recorded in the package manifest and result record; the
mathematical coverage is then recomputed. From the repository root, run
\begin{verbatim}
python3 certificates/lower_bound_4p614153/verify.py \
  --output /tmp/n17-replay --upstream --direct
python3 certificates/lower_bound_4p614153/test_verify.py
python3 certificates/lower_bound_4p614153/check_diagnostics.py
\end{verbatim}
The output directory must not already exist. Requirements are Python 3.10+,
a C++17 compiler with OpenMP, and Boost headers. Numerical optimization and
network access are not part of verification.

The current catalogue requires 27,918,671 center slabs. The two accumulation
modes, a segment tree and a direct prefix calculation, must agree on every
interval minimum. They share the geometric partition and do not constitute
two independent geometric proofs. With \texttt{--upstream}, the command also
reconstructs and checks all 2,925 R012 intervals from the pinned source
in~\cite{Guzhou}; it does not run that source's Python sweep. Earlier research
runs of the source-distinct Python implementation are documented separately
in the result record.

Geometric comparisons use unbounded integers or exact rational arithmetic.
Integer weight accumulators are signed 64-bit values; validation bounds the
total nonnegative integer weight by $10^{12}$, below the accumulator limits.
The checker requires a complete output for every catalogue row and verifies
the counting inequality using the recomputed minima.

The recentering check verifies every inequality~\eqref{eq:scalar}, the bound
on $\sqrt2$, positivity of the new parameters, and nonempty legal-center
domains. Full exact values of its minimum slack and the cap in
Appendix~\ref{app:cap} are retained in \path{result.json} and in the replay's
\path{sharpening.result.json}. The support and catalogue were found by numerical
search; their validity depends only on this exact replay and the reductions
proved in the main text.

\subsection{Trigger certificate}

The directory \path{certificates/intersection_trigger_29/} contains
\path{intersection-atom.json}, \path{clique-29-poses-with-weights.csv}, the
pinned source data, and an exact Python checker. Run
\begin{verbatim}
python3 certificates/intersection_trigger_29/verify.py
python3 certificates/intersection_trigger_29/test_verify.py
\end{verbatim}
Only Python 3.10+ and its standard library are needed. The checker validates
rational unit rotations, containment of the poses in the side-$4.67$ box,
every shared witness and inward margin, and all fractional weights. It
computes the 29-pose load across 1,475 arrangement slabs, including a direct
recount of a maximizing point. Source hashes identify the supplied artifacts;
the separation proof checks the selected poses directly and assumes no load
bound for the full source family.

\section{Historical and supplementary results}\label{app:history}

\subsection{Selected lower-bound history}

Table~\ref{tab:lineage} records the results that shaped this work and nearby
benchmarks. Classical unavoidable-point arguments lead to the integral
counting rule; fractional measures and restricted parent-center domains
provide the successive relaxations used here. Other geometric approaches
include the work of Stromquist, Nagamochi, and
MacIver~\cite{Stromquist2003,Nagamochi2005,MacIver}. Detailed attribution and
date evidence are retained in \path{CONTRIBUTORS.md} and
\path{paper/HISTORY-DATES.md}.

\begin{table}[!htbp]
\centering
\caption{Selected $n=17$ lower bounds, ordered by bound rather than date.
Dates identify publication or retained records, not necessarily discovery;
Green's is the reported communication year. Rounded values marked $\approx$
are only summaries of their source bounds.}
\label{tab:lineage}
\small
\begin{tabular}{@{}L{0.16\textwidth}L{0.23\textwidth}L{0.20\textwidth}L{0.33\textwidth}@{}}
\toprule
Date & Contributor / source & Bound & Contribution\\
\midrule
2000 & Green, via Friedman~\cite{Friedman2009} & $\approx4.445208$ & Sixteen unavoidable points.\\
2026-07-18 & Brandwijk~\cite{Brandwijk} & $4.45$ & Exact point certificate.\\
2026-08-10; 2026-08-11 & Mira~\cite{MiraRepo} & $4.450837$; $4.468292$ & Exact pose subdivision; later triangle witnesses.\\
2026-08-11 & Fort~\cite{Fort} & $4.456575$ & Independent exact pose subdivision.\\
2026-08-06 & Burns~\cite{Burns} & $4.4811$ & Fractional atomic measure.\\
2026-08-21 & Massaccesi~\cite{Massaccesi} & $4.5058$ & LP-designed weighted support.\\
2026-08-16 & anabologyco-maker~\cite{Anabology} & $4.5705$ & Weighted certificate with a Lean finite-check layer.\\
2026-09-04 & Levy~\cite{Levy} & $4.59$ & Reusable exact event sweeps and coverage checks.\\
2026-09-07 & Mira continuation~\cite{MiraRepo} & $\approx4.607029$; $\approx4.613029$ & Support and angle refinement; exact dilation.\\
2026-09-19 & Guzhou0806 / R012~\cite{Guzhou} & $\approx4.613046$ & Parent-angle catalogue and legal-center coverage.\\
2026-09-19 & This work & $\approx4.614154$ & New measure and catalogue; analytic recentering.\\
\bottomrule
\end{tabular}
\end{table}

R012's exact constant is $461300/99999$. Intermediate search reports reached
$922600/199979$ with its measure fixed, and then $230650/49993$ by reweighting
the same support. Their separate certificates are not retained here. The
published replay covers R012, the final base certificate, and the recentering
corollary.

An exact retained diagnostic explains one limit of changing only the core
catalogue. At parent side $99989/100000$ in the side-$4.613$ container, a legal
parent has R012 mass $0.998982$, below the mass required by its common-charge
counting inequality. No inner core can have more mass than the parent. This
obstructs that fixed measure at target $461300/99989$, rather than bounding
$s(17)$ or the general method. The final certificate changes the measure by
reweighting inherited sites and adding wall-adjusted support.

\subsection{A limit on reusing the fixed core domains}\label{app:cap}

The recentering argument admits a simple upper bound on what the unchanged
core shapes and certified domains alone can prove. For an axis-aligned parent,
write $W_j=B_jf(t_j)$. Lemma~\ref{lem:recenter} requires
\[
  a>W_j,\qquad H<2a+2d_j-W_j.
\]
These inequalities follow already by considering a corner parent center.
Thus even if the choice of entry depends on that center, some entry must
satisfy them. For $a>W_j$,
\[
  \frac Ha<2+\frac{2d_j-W_j}{a}
  \le\max\left\{2,\,1+\frac{2d_j}{W_j}\right\}.
\]
Consequently any such reuse argument has
\[
  \frac Ha\le C_*:=\max_j\max\left\{2,\,1+\frac{2d_j}{W_j}\right\},
  \qquad C_*\approx4.614153538422734.
\]
The exact rational maximum is checked in the scalar replay. This is a limit
on those unchanged core/domain statements, not an upper bound on $s(17)$ or
on certificates using new shapes, domains, measures, or charges.

\subsection{Additional trigger activations}

The trigger checker also verifies 71 unit-square activations from the pinned
source, of which 68 remain valid after shrinking to side $499/500$. The
68 retained weights sum to $96836866527/79999999999$. A separate exact
arrangement calculation on the 71 unit squares gives maximum pointwise load
$73959042636/79999999999<1$. The integration argument used in
Proposition~\ref{prop:separation} therefore verifies the earlier, weaker
mass lower bound $96836866527/79999999999>1.21$ on the 68 cores. The 29-core
result in the main text is stronger. Neither calculation assumes a load
bound for all 616 poses in the larger source family.

\clearpage
\begin{thebibliography}{99}

\bibitem{Friedman2009}
E. Friedman,
\newblock \emph{Packing Unit Squares in Squares: A Survey and New Results},
\newblock Electron. J. Combin. Dynamic Survey DS7 (2009).
\newblock \url{https://erich-friedman.github.io/papers/squares/squares.html}.

\bibitem{Stromquist2003}
W. Stromquist,
\newblock Packing 10 or 11 unit squares in a square,
\newblock Electron. J. Combin. 10(1) (2003), R8.
\newblock \url{https://doi.org/10.37236/1701}.

\bibitem{Nagamochi2005}
H. Nagamochi,
\newblock Packing unit squares in a rectangle,
\newblock Electron. J. Combin. 12(1) (2005), R37.
\newblock \url{https://doi.org/10.37236/1934}.

\bibitem{Ellsworth}
D. Ellsworth,
\newblock Squares in squares, high-precision online packing catalogue.
\newblock \url{https://kingbird.myphotos.cc/packing/squares_in_squares.html}.

\bibitem{MiraRepo}
Mira,
\newblock 17squares: exact computer-assisted certificates and research archive.
\newblock \url{https://github.com/Mira-acc/17squares}.

\bibitem{Fort}
S. Fort,
\newblock Packing 17 unit squares: exact lower bound $s(17)>4.456575$.
\newblock \url{https://github.com/stanislavfort/17squares}.

\bibitem{Burns}
S. Burns,
\newblock Proposing a better lower bound for $n=17$ square packing, 2026.
\newblock \url{https://sam-burns.com/posts/proposing-better-lower-bound-for-n17-square-packing/}.

\bibitem{Brandwijk}
K. Brandwijk,
\newblock Exact $n=17$ lower-bound certificate at $89/20$, 2026.
\newblock Zenodo, DOI 10.5281/zenodo.21422426.
\newblock \url{https://doi.org/10.5281/zenodo.21422426}.

\bibitem{Massaccesi}
G. Massaccesi,
\newblock Another better lower bound for $n=17$ square packing, 2026.
\newblock \url{https://gus-massa.blogspot.com/2026/08/another-better-lower-bound-for-n17.html}.
\newblock See also \emph{Linear Programing for Square Packing},
\url{https://gus-massa.blogspot.com/2026/08/linear-programing-for-square-packing.html}.

\bibitem{Anabology}
anabologyco-maker,
\newblock Compressed exact certificate candidate for packing 17 unit squares: side at least $4.5705$.
\newblock \url{https://github.com/anabologyco-maker/square17-lower-bound}.

\bibitem{Levy}
J. Levy,
\newblock The squares project: weighted fractional certificates, exact event sweeps, and frontier audit.
\newblock \url{https://github.com/jlevy/squares}.

\bibitem{Guzhou}
Guzhou0806 / N17 project,
\newblock R012 parent-angle lower-bound certificate.
\newblock \url{https://github.com/Guzhou0806/n17-square-packing/tree/931a0dfd64302e277057006e99388fe5c00b7f53/certificates/R012}.

\bibitem{MacIver}
D. R. MacIver,
\newblock An improved lower bound for packing seventeen unit squares in a square, 2026.
\newblock \url{https://github.com/DRMacIver/square-packing-research/blob/9e2cd597047e040a63b7dcd103e79962cfc2d781/papers/s17-lower-bound/paper.pdf}.

\end{thebibliography}

\end{document}
